% Сборка из корня репозитория: latexmk exam/theoretical-minimum.tex
\documentclass[11pt,a4paper]{article}
\newcommand{\examtitle}{Теоретический минимум}
\usepackage{iftex}
\RequirePDFTeX
\usepackage{cmap}
\usepackage[T2A]{fontenc}
\usepackage[utf8]{inputenc}
\usepackage[english,russian]{babel}
\usepackage{PTSerif}
\usepackage{amsmath,amssymb,microtype}
\usepackage[a4paper,left=23mm,right=23mm,top=20mm,bottom=21mm]{geometry}
\usepackage{xcolor,enumitem,fancyhdr,hyperref}

\definecolor{accent}{HTML}{187AA5}
\definecolor{ink}{HTML}{302B28}
\definecolor{muted}{HTML}{6D6560}
\definecolor{tint}{HTML}{EDF4F6}
\definecolor{rulecolor}{HTML}{A9CADA}
\hypersetup{unicode=true,colorlinks=true,urlcolor=accent,linkcolor=accent,
  pdftitle={\examtitle},pdfauthor={Чиняев Владимир Николаевич}}
\color{ink}
\setlength{\parindent}{0pt}
\setlength{\parskip}{3pt}
\linespread{1.08}
\setlist[enumerate]{leftmargin=*,label=\textcolor{accent}{\arabic*.},
  widest=99,itemsep=3pt,topsep=4pt,parsep=0pt}
\setlength{\emergencystretch}{1em}
\pagestyle{fancy}
\fancyhf{}
\renewcommand{\headrulewidth}{0pt}
\renewcommand{\footrulewidth}{0.4pt}
\renewcommand{\footrule}{\hbox to\headwidth{\color{rulecolor}\leaders\hrule height\footrulewidth\hfill}}
\fancyfoot[L]{\footnotesize\color{muted}Введение в маневрирование КА\quad /\quad \examtitle}
\fancyfoot[R]{\footnotesize\color{muted}\thepage}

\newcommand{\examheading}[1]{%
  {\small\color{accent}ВВЕДЕНИЕ В МАНЕВРИРОВАНИЕ КА}\par
  \vspace{3pt}
  {\huge\examtitle}\par
  {\small\color{muted}#1}\par
  \vspace{4pt}
  {\color{rulecolor}\hrule}
  \vspace{5pt}}

\newcommand{\examnote}[1]{%
  \par\smallskip\noindent
  \begingroup
  \setlength{\fboxsep}{7pt}%
  \setlength{\fboxrule}{0.4pt}%
  \fcolorbox{rulecolor}{tint}{%
    \begin{minipage}{\dimexpr\linewidth-2\fboxsep-2\fboxrule\relax}
      \small #1\par
    \end{minipage}}%
  \endgroup\par\smallskip}

\newcommand{\topicheading}[1]{%
  \par\vspace{5pt}{\large\bfseries\color{accent}#1}\par\nobreak\smallskip}


\begin{document}
\examheading{Допуск к дифференцированному зачёту \enspace / \enspace 2026--2027 учебный год}

\examnote{\textbf{Правила сдачи теоретического минимума.}
Задаются три вопроса из списка. Для допуска к зачёту нужно правильно и быстро
ответить на все три. Теоретический минимум баллов не даёт. Если минимум не сдан,
итоговая оценка за курс~--- 2 независимо от баллов за домашние задания.}

\topicheading{Вопросы}

\begin{enumerate}
\item {\raggedright Перечислите шесть классических кеплеровых элементов.
Укажите единицы \mbox{измерения. Изобразите} их на одном рисунке.\par}

\item Запишите радиусы перицентра и апоцентра эллиптической орбиты через
большую полуось и эксцентриситет.

\item Как направлены радиальная, трансверсальная и нормальная оси
орбитальной системы координат?

\item Какую особенность формы Земли учитывает возмущение $J_2$? Какие кеплеровы
элементы испытывают вековые изменения в первом приближении по $J_2$?

\item Что происходит с положением и скоростью КА при мгновенном импульсе?
Запишите связь скоростей до и после импульса.

\item Какому геометрическому условию должны удовлетворять две заданные орбиты,
чтобы между ними был возможен одноимпульсный переход?

\item КА на круговой орбите получил малый импульс, сонаправленный со скоростью.
В какой апсиде новой орбиты приложен импульс? Как изменился период обращения?

\item Какая компонента импульса поворачивает плоскость околокруговой орбиты
в первом приближении?

\item Чему равен модуль базис-вектора Лоудена в точке ненулевого импульса
при проверке оптимальности? Как направлен этот вектор относительно импульса?

\item Какие конечные условия задают в задаче перехода?

\item Какие конечные условия задают в задаче встречи?

\item Что задано и что требуется найти в задаче Ламберта?

\item Какую опорную орбиту предполагает модель Клохесси--Уилтшира? Каким должно
быть расстояние между КА и целью по сравнению с радиусом этой орбиты?

\item Как найти тяговое ускорение по силе тяги и текущей массе КА?

\item Как ориентирована плоскость сближения относительно вектора относительной
скорости в момент наибольшего сближения?

\item Чему должны равняться радиальная и трансверсальная компоненты скорости
при выходе на круговую орбиту радиуса $r$ в центральном поле?

\item Какой вектор оценивают по орбитальным измерениям при реконструкции
выполненного импульсного манёвра?
\end{enumerate}
\end{document}
